\begin{ourthm}
\label{T: Top Las}
Let $\alpha$ and $\gamma$ be two weak compositions. 
\begin{enumerate}
\item[(a)]
The polynomial $\topLas_\alpha$ has leading monomial $x^{\rajcode(\alpha)}$.
\item[(b)] We have
$\topLas_\alpha$ is a scalar multiple of $\topLas_{\gamma}$
if and only if $\rajcode(\alpha) = \rajcode(\gamma)$.
\item[(c)] 
We say $\alpha$ is snowy if its
positive entries are distinct.
If $\alpha$ is snowy, then $x^{\rajcode(\alpha)}$ has coefficient 1 in $\topLas_\alpha$.
Moreover, 
there exists exactly one snowy weak composition $\gamma'$ 
such that $\rajcode(\gamma) = \rajcode(\gamma')$.
\end{enumerate}     
\end{ourthm}

Our definition of $\rajcode(\cdot)$
on weak compositions is diagrammatic. 
Given a diagram~$D$, 
we define a combinatorial
construction called the snow diagram 
that augments and decorates $D$.
Let $\rajcode(D)$
be the weight of the snow diagram.
Every weak composition $\alpha$ is naturally
associated with a diagram called
the key diagram $D(\alpha)$
(see Subsection~\ref{subsec.diagrams}).
Then we define $\rajcode(\alpha):=\rajcode(D(\alpha))$.

Snow diagrams unify the computation of leading monomials in $\topGro_w$ and $\topLas_\alpha$.
Each permutation $w$ is also associated
with a diagram called the Rothe diagram
$RD(w)$.
In~\S\ref{S: permutations},
we show $\rajcode(w) = \rajcode(RD(w))$.
In other words, we give a diagrammatic way to compute $\rajcode(w)$.

Finally, let $\widehat{V}_n :=  \Q\textrm{-span}\{\topGro_w: w \in S_n\}$ and $\widehat{V} :=  
\bigcup_{n \geq 1} \widehat{V}_n$.
In Proposition~\ref{P: Filtered Algebra},
we show $\widehat{V}$
is a filtered algebra. 
Theorem~\ref{thm:psw} can be used to construct a basis of $\widehat{V}_n$
and~$\widehat{V}$ 
consisting of $\topGro_w$.
In particular, 
the dimension of $\widehat{V}_n$
is $B_n$, the $n\textsuperscript{th}$ Bell number.
In~\S\ref{S: Space},
we use Theorem~\ref{T: Top Las} to construct another basis consisting of $\topLas_\alpha$.
This basis allows us to compute the Hilbert series of $\widehat{V}_n$ 
and $\widehat{V}$
involving a $q$-analogue of $B_n$.

The rest of the paper is organized as follows. 
In \S\ref{S: Background}, we provide necessary background information and notation.
In \S\ref{S: Cloud and snow}, we construct a snow diagram from any diagram 
and define statistics $\rajcode(\cdot)$ and $ \raj(\cdot)$ on all diagrams.
In \S\ref{S: weak composition}, we  prove Theorem~\ref{T: Top Las}.
In~\S\ref{S: permutations}, we show the statistics $\rajcode(\cdot)$ and $\raj(\cdot)$ 
on a Rothe diagram are equivalent to that defined in~\cite{PSW}.
We also relate the snow diagram to two classical constructions: 
Schensted insertion and the shadow diagram.
In \S\ref{S: Space}, we derive the Hilbert series 
of $\widehat{V}_n$ and $\widehat{V}$.
In \S\ref{sec.open}, we present several open problems and future directions.

\section{Background}
\label{S: Background}
\subsection{Polynomials}
We provide necessary background
for Grothendieck polynomials and Lascoux polynomials.
Then we introduce
$\topGro_w$ and $\topLas_\alpha$ 
which span the spaces $\widehat{V}_n$ and $\widehat{V}$. 

The \definition{Grothendieck polynomials} $\fGb_w \in 
\mathbb{Z}_{\geqslant 0}[x_1, x_2, \dots][\beta]$
were recursively defined by Lascoux and Sch{\"u}tzenberger~\cite{LS:Groth}. 
Let $\partial_i(\cdot)$ be the divided difference operators acting 
on the polynomial ring.
For each $i$,
define $\partial_i(f) := \dfrac{f - s_i f}{x_i-x_{i+1}}$,
where $s_i$ is the operator that swaps $x_i$ and $x_{i+1}$.
Then for $w\in S_n$, 
\begin{align*}
\fGb_w &:=
\begin{cases}
x_1^{n-1}x_2^{n-2}\cdots x_{n-1} &\text{if $w$ is $[n, n-1, \dots, 1]$ in one-line notation,} \\
\partial_i( (1 + \beta x_{i+1}) \fGb_{ws_i} ) &\text{if $w(i)<w(i+1)$.}
\end{cases}
\end{align*}

Let $S_+$ be the set of permutations
of $\{1, 2, \dots\}$
such that only finitely many numbers are permuted.
Take $w \in S_+$ and assume $w$ only permutes numbers in $[n]$.
Let $w' \in S_n$ be the restriction of $w$ to $[n]$
and define $\fGb_w$ as $\fGb_{w'}$.
It is shown in~\cite{LS:Groth} that $\fGb_w$ is well-defined.

A \definition{weak composition}
is an infinite sequence of non-negative integers with finitely many positive entries. 
Let $C_+$ be the set of weak compositions. 
For $\alpha \in C_+$, 
we use $\alpha_i$ to denote its $i$\textsuperscript{th} entry,
and write $\alpha = (\alpha_1, \alpha_2, \dots, \alpha_n)$ 
where $\alpha_n$ is the last positive entry. 
We use $x^\alpha$ to denote the monomial
$x_1^{\alpha_1}x_2^{\alpha_2}\cdots x_n^{\alpha_n}$ and
$|\alpha| = \sum_{i\geqslant 1}^n\alpha_i$.
The \definition{Lascoux polynomials} 
$\fLb_\alpha$, indexed by weak compositions,
are in $\mathbb{Z}_{\geqslant 0}[x_1, x_2, \dots][\beta]$. 
By~\cite{Las}, they are defined recursively
\begin{align*}
\fLb_\alpha = \begin{cases}
x^\alpha & \text{if $\alpha$ is weakly decreasing,} \\
\pi_i ((1 + \beta x_{i+1}) \fLb_{s_i\alpha} ) &\text{if $\alpha_i<\alpha_{i+1}$,}
\end{cases}
\end{align*}
where $\pi_i$ is the operator $\pi_i(f) := \partial_i(x_i f)$.


We say a pair $(i,j)$ is an \definition{inversion}
of $w \in S_n$ if $i < j$ and $w(i) > w(j)$. 
Let $\Inv(w)$ be the set of all inversions in $w$
and let $\inv(w) = |\Inv(w)|$.
Then we may view $\fGb_w$ as a polynomial in $\beta$, where
$$
[\beta^d] \fGb_w \quad := \quad \textrm{coefficient of $\beta^d$ in $\fGb_w$}
$$
is a homogeneous polynomial in the $x$-variables with degree 
$\inv(w) + d$ in $\Z_{\geq 0}[x_1,x_2,\dots]$.
The \definition{Schubert polynomial} $\fS_w := [\beta^0]\fGb_w$.
Similarly, viewing $\fLb_\alpha$ as a polynomial of $\beta$,
$[\beta^d] \fLb_\alpha$ is a homogeneous polynomial with degree
$|\alpha| + d$ in $\Z_{\geq 0}[x_1,x_2,\dots]$.
The \definition{key polynomial} $\kappa_\alpha := [\beta^0]\fLb_\alpha$.
The representation theoretic, geometric and combinatorial perspectives of Schubert polynomials and key polynomials
are well-studied~\cite{De,LS1,BJS}.

Define $V_n :=  \Q\textrm{-span}\{\fS_w: w \in S_n\}$ 
and $V :=  \Q\textrm{-span}\{\fS_w: w \in S_+\} = \bigcup_{n \geq 1} V_n$.
In fact, $V = \mathbb{Q}[x_1, x_2, \dots]$. 
By the increasing sequence 
$V_1 \subset V_2 \subset \cdots \subset V$,
$V$ has the structure of a filtered algebra. 


In this paper, we are interested in the top-degree 
components of $\fGb_w$ and $\fLb_\alpha$. 
For a polynomial $f \in \mathbb{Q}[x_1, x_2, \dots][\beta]$,
let $\widehat{f} = [\beta^d](f)$ where $d$ is the largest such that 
$[\beta^d](f) \neq 0$.
The \definition{Castelnuovo--Mumford polynomial} of $w \in S_+$
is defined as~$\topGro_w$.
The \definition{top Lascoux polynomial} of $\alpha \in C_+$
is defined as~$\topLas_\alpha$. In appendix \S\ref{appendix}, we list some Grothendieck polynomials and Lascoux polynomials.
Pechenik, Speyer and Weigandt~\cite{PSW} first study $\topGro_w$.
To the best of the authors knowledge, 
$\topLas_\alpha$ has not been studied previously. 

Now consider the \definition{tail lexicographic order} 
on monomials in the $x$-variables.
We say a monomial $x^\alpha$ is larger than $x^\gamma$
if there exists $k$ such that $\alpha_k > \gamma_k$ and 
$\alpha_j = \gamma_j$ for all $j > k$.
The \definition{leading monomial} of $f \in \Q[x_1, x_2, \dots]$
is the largest monomial in $f$. 
Among the four homogeneous
polynomials above,
three of them have combinatorial rules for their leading terms:
\begin{enumerate}
\item~\cite{BJS} The leading monomial of $\fS_w$ with $w \in S_n$
is $x^{\invcode(w)}$,
where $$\invcode(w)_i = |\{j: (i,j) \in \Inv(w)\}|.$$
\item~\cite{LS2} The leading monomial of $\kappa_\alpha$ is $x^\alpha$.
\item~\cite{PSW} The leading monomial of $\topGro_w$
is $x^{\rajcode(w)}$ defined as follows.
\end{enumerate}

\begin{defn}\cite{PSW}
\label{D: raj on permutations}
Let $\textup{LIS}^w(q)$ be the 
length of the longest increasing subsequence 
of $w\in S_n$ that starts with $q$.
The $\rajcode(w)$ for $w\in S_n$ is a weak composition
where 
$\rajcode(w)_r := n + 1 - r - \textup{LIS}^w(w(r))$ for $r \in [n]$
and $0$ if $r > n$.
Then $\raj(w) := |\rajcode(w)|$.
\end{defn}


\begin{exa}
\label{E: rajcode permutation}
Consider $w = 3721564 \in S_7$. We have $\textup{LIS}^w(2) = 3$, 
so $\rajcode(w)_3 = 7 + 1 - 3 - 3 = 2$.
All together,
we get $\rajcode(w) = (4,5,2,1,1,1)$ and $\raj(w) = 14$.
\end{exa}

We will define $\rajcode(\cdot)$ on $C_+$
and show the leading monomial of $\fLb_\alpha$ is $x^{\rajcode(\alpha)}$
in \S\ref{S: weak composition}.

A connection between $\fGb_w$ and $\fLb_\alpha$
is established by Shimozono and Yu~\cite{SY}.
To describe this connection, 
we need the following notion.
\begin{defn}
Let $f, f_1, f_2, \dots$ be polynomials in 
$\mathbb{Z}_{\geq 0}[x_1, x_2,\dots]$.
We say $f$ \definition{expands positively} 
into $\{f_1,f_2, \dots\}$ if there exist 
$c_1,c_2, \dots \in \mathbb{Z}_{\geq 0}$
such that $f = \sum_{i} c_if_i$.

Now assume $f, f_1, f_2, \dots$ are polynomials 
in $\mathbb{Z}_{ \geq 0}[\beta][x_1, x_2,\dots]$.
We say $f$ \definition{expands positively} 
into $\{f_1,f_2, \dots\}$ if there exist 
$g_1,g_2, \dots \in \mathbb{Z}_{\geq 0}[\beta]$ 
such that $f = \sum_{i} g_if_i$.
\end{defn}

\begin{ourthm}[\cite{SY}]
\label{T: Gro to Las}
For $w \in S_+$, $\fGb_w$ expands positively
into $\{\fLb_\alpha: \alpha \in C_+\}$. 
\end{ourthm}
 
This result implies
$\topGro_w$ also expands
positively into
$\topLas_\alpha$ by the following lemma whose proof
is sufficiently elementary. 

\begin{lemma}
\label{L: Expands positively}
Let $f, f_1, f_2, \dots$ 
in $\mathbb{Z}_{\geq 0}[\beta][x_1, x_2,\dots]$.
If $f$ expands positively 
into $\{f_1, f_2, \dots\}$,
then $\widehat{f}$ expands positively into $\widehat{f_1}, \widehat{f_2}, \dots $.
\end{lemma}

\begin{cor}
\label{C: top Gro to top Las}
For $w \in S_+$,
$\topGro_w$ expands positively into $\{\topLas_\alpha: \alpha \in C_+\}$.
\end{cor}

Define $\widehat{V}_n :=  \Q\textrm{-span}\{\topGro_w: w \in S_n\}$ and $\widehat{V} :=  \Q\textrm{-span}\{\topGro_w: w \in S_+\} = \bigcup_{n \geq 1} \widehat{V}_n$.
By work of Lascoux, Sch{\"u}tzenberger~\cite{LS:Groth} and Brion~\cite{Bri},
the product $\fGb_u \fGb_v$ with $u \in S_m$ and $v \in S_n$
expands positively into $\fGb_w$ with $w \in S_{m + n}$.
By Lemma~\ref{L: Expands positively},
$\topGro_u \topGro_v$ with $u \in S_m$ and $v \in S_n$ expands positively into $\topGro_w$ 
with $w \in S_{m + n}$.
Finally, we conclude the following.
\begin{prop}
\label{P: Filtered Algebra}
The space $\widehat{V}$ is a filtered 
algebra with respect to the filtration
$\widehat{V}_1 \subset \widehat{V}_2 \subset \cdots \subset \widehat{V}$.
\end{prop}


\subsection{Diagrams}\label{subsec.diagrams}
A \definition{diagram} is a finite subset 
of $\mathbb{Z}_{>0} \times \mathbb{Z}_{>0}$.
We represent a diagram by putting a cell at row $r$
and column $c$ for each $(r,c)$ in the diagram.
The leftmost column (resp. topmost row) is called column $1$ (resp. row $1$).
The \definition{weight} of a diagram~$D$, 
denoted as $\wt(D)$, 
is a weak composition 
whose $i$\textsuperscript{th}  entry is the number of boxes 
in its row $i$.
We recall two classical families of diagrams. 

Each weak composition $\alpha$
is associated with a diagram called 
the \definition{key diagram}, 
denoted as $D(\alpha)$.
It is the unique left-justified diagram 
with weight $\alpha$.
One important key diagram we will use later is 
$\Stair_n := D((n-1, n-2, \cdots, 1))$.

\begin{exa}
The following are two examples of key diagrams.
For clarity, we put an ``$i$'' 
on the left of row $i$
and put a small dot in each cell.
\begin{align*}
D(0,2,1) = 
\raisebox{0.45cm}{
\begin{ytableau}
\none[1]\cr
\none[2] & \cdot & \cdot \cr
\none[3] & \cdot
\end{ytableau}}
\quad , \quad \quad \quad
\Stair_4 = 
\raisebox{0.45cm}{
\begin{ytableau}
\none[1] & \cdot & \cdot & \cdot \cr
\none[2] & \cdot & \cdot \cr
\none[3] & \cdot
\end{ytableau}} 
\end{align*}
\end{exa}

Each permutation $w$ 
is associated with
the \definition{Rothe diagram} 
$RD(w):=\{ (r, w(r'): (r,r') \in \Inv(w)\}$.

\begin{exa}
Let $w = 41532 \in S_5$.
Then $$\Inv(w) = \{ (1,2), (1,4), (1,5), 
(3,4), (3,5), (4,5)\}\,.$$
The Rothe diagram is depicted as follows.
$$
RD(w) = 
\raisebox{1cm}{
\begin{ytableau}
\none[1] & \cdot & \cdot & \cdot \cr
\none[2] & \none & \none & \none \cr
\none[3] & \none & \cdot & \cdot \cr
\none[4] & \none & \cdot & \none \cr
\none[5] & \none & \none & \none 
\end{ytableau}}
$$
\end{exa}

\subsection{\texorpdfstring{$K$}{TEXT}-Kohnert diagrams}
We recall a combinatorial formula for Lascoux polynomials. 
To simplify our description, 
we introduce the following definition.
\begin{defn}
A \definition{labeled diagram} 
is a diagram where each cell can be labeled 
by a symbol.
The \definition{underlying diagram} of a labeled diagram
is the diagram obtained by ignoring all labels. 
The weight of a labeled diagram $D$, denoted as $\wt(D)$,
is just the weight of its underlying diagram.
\end{defn}

Then a \definition{ghost diagram} is a labeled diagram 
where cells can be labeled by $\X$.
We call cells labeled by $\X$ as ``\definition{ghosts}''. 
For a ghost diagram $D$, its \definition{excess}, denoted as $\ex(D)$,
is the number of ghosts in $D$.
Next, we define a move on ghost diagrams.

\begin{defn}[\cite{RossY}]
A \definition{$K$-Kohnert move} is defined on a ghost diagram $D$. 

We pick a cell $(r,c)$ and move it up, 
subject to the following requirements.
\begin{itemize}
\item The cell $(r,c)$ must be the rightmost cell in row $r$.
\item The cell $(r,c)$ is not a ghost.
\item The cell $(r,c)$ is moved to the lowest empty spot above it.
\item The cell $(r,c)$ may jump over other cells but cannot jump over any ghosts.
\end{itemize}
After the move, 
we may or may not leave a ghost at $(r,c)$. 
When we leave a ghost, 
we refer this move as a \definition{ghost move}.
\end{defn}

For a weak composition $\alpha$,
a ghost diagram is called a \definition{$K$-Kohnert diagram}
of $\alpha$ if it can be obtained from $D(\alpha)$
by $K$-Kohnert moves. 
Let $\KKD(\alpha)$ be the set of all $K$-Kohnert diagrams
of $\alpha$.
As proved in~\cite{PY}, 
$K$-Kohnert diagrams give a formula for Lascoux polynomials. 
This rule was first conjectured by Ross and Yong~\cite{RossY}.
Notice that our convention is different from~\cite{PY}:
row $1$ is the top most row in this paper while it is the 
bottom most row in~\cite{PY}. 

\begin{ourthm}[\cite{PY}]\label{thm: k.kohnert}
Let $\alpha$ be a weak composition.
Then we have
$$
\fLb_\alpha = \sum_{D \in \KKD(\alpha)}x^{\wt(D)} \beta^{\ex(D)}.
$$
\end{ourthm}

\begin{exa} \label{eg_021} Let $\alpha = (0,2,1)$, then 
$KKD(\alpha)$ consists of the following: 
\begin{align*}
\raisebox{1.1cm}{
\begin{ytableau}
\none[1] \cr
\none[2] & \cdot & \cdot \cr
\none[3] & \cdot 
\end{ytableau}}\:,
\raisebox{1.1cm}{
\begin{ytableau}
\none[1] &\none & \cdot \cr
\none[2] & \cdot \cr
\none[3] & \cdot 
\end{ytableau}}\:,
\raisebox{1.1cm}{
\begin{ytableau}
\none[1] & \cdot & \cdot \cr
\none[2]  \cr
\none[3] & \cdot 
\end{ytableau}}\:,
\raisebox{1.1cm}{
\begin{ytableau}
\none[1] & \cdot \cr
\none[2] & \cdot & \cdot \cr
\none[3] 
\end{ytableau}}\:,
\raisebox{1.1cm}{
\begin{ytableau}
\none[1] & \cdot & \cdot \cr
\none[2] & \cdot  \cr
\none[3] 
\end{ytableau}}\:,
\end{align*}
\begin{align*}
\raisebox{1.1cm}{
\begin{ytableau}
\none[1] & \cdot \cr
\none[2] & \cdot & \cdot \cr
\none[3] & \X 
\end{ytableau}}\:,   
\raisebox{1.1cm}{
\begin{ytableau}
\none[1] &\none & \cdot \cr
\none[2] & \cdot & \X \cr
\none[3] & \cdot 
\end{ytableau}}\:,
\raisebox{1.1cm}{
\begin{ytableau}
\none[1] & \cdot & \cdot \cr
\none[2] & \X \cr
\none[3] & \cdot 
\end{ytableau}}\:,
\raisebox{1.1cm}{
\begin{ytableau}
\none[1] & \cdot & \cdot \cr
\none[2] & \cdot \cr
\none[3] & \X
\end{ytableau}}\:,
\raisebox{1.1cm}{
\begin{ytableau}
\none[1] & \cdot & \cdot \cr
\none[2] & \cdot & \X \cr
\none[3] 
\end{ytableau}}\:,
\raisebox{1.1cm}{
\begin{ytableau}
\none[1] & \cdot & \cdot \cr
\none[2] & \cdot & \X\cr
\none[3] & \X
\end{ytableau}}\:.   
\end{align*}

By the rule above, we have
\begin{equation*}
\begin{split}
\fLb_{\alpha} = & \: x_2^2x_3 + x_1x_2x_3 + x_1^2x_3 
+ x_1x_2^2 + x_1^2x_2\\
+ & \: \beta(x_1x_2^2x_3 + x_1x_2^2x_3 + x_1^2x_2x_3 + x_1^2x_2x_3 + x_1^2x_2^2) + \beta^2 x_1^2x_2^2x_3.
\end{split}
\end{equation*}
\end{exa}


\section{Snow diagrams}
\label{S: Cloud and snow}
We associate each diagram with a labeled diagram
called the \definition{snow diagram} which allows us to define two statistics
on diagrams. 
For each diagram $D$, 
we describe the following algorithm that outputs $\snow(D)$.
Cells in $\snow(D)$ can be labeled by $\bullet$ or~$\snowflake$.


\begin{enumerate}
\item[-] Iterate 
through rows of $D$ from bottom to top. 
\item[-] In each row $r$ of $D$, find the rightmost cell $(r,c)$ 
with no $\bullet$ in column $c$.
If such an $(r,c)$ exists,
label it by $\bullet$
and put a cell labeled by $\snowflake$ in $(r', c)$ for $r' \in [r-1]$ and $(r', c) \notin D$.
\end{enumerate}

We call cells labeled by $\bullet$ \definition{dark clouds}
and cells labeled by $\snowflake$  \definition{snowflakes}.

\begin{exa}\label{exa.snow}
The following is a diagram together with its snow diagram. 
\begin{align*}
D = 
\raisebox{0.95cm}{
\begin{ytableau}
\none[1] & \none& \none & \cdot \cr
\none[2] & \cdot& \cdot & \none \cr
\none[3] & \none& \none & \cdot \cr
\none[4] & \none& \none & \none \cr
\none[5] & \cdot& \cdot \cr
\end{ytableau}}\,,\quad
\snow(D) = 
\raisebox{0.95cm}{
\begin{ytableau}
\none[1] & \snowflake& \snowflake & \cdot \cr
\none[2] & \bullet& \cdot & \snowflake \cr
\none[3] & \none& \snowflake & \bullet \cr
\none[4] & \none& \snowflake & \none \cr
\none[5] & \cdot& \bullet \cr
\end{ytableau}}\quad.
\end{align*}
\end{exa}

The positions of dark clouds will be important,
so we make the following definition.

\begin{defn}
The \definition{dark cloud diagram} of a diagram $D$,
$\dark(D)$, 
is the set of cells $(r,c)$ that are dark clouds in $\snow(D)$.
\end{defn}

\begin{exa}
In Example~\ref{exa.snow},
$\dark(D) = \{(2,1), (3,3), (5,2) \}$.
\end{exa}

A diagram is a 
\definition{non-attacking rook diagram}
if it has at most one cell in each row or column.
Let $\Rook_+$ be the family
of all non-attacking rook diagrams.  


\begin{rem}
\label{R: Basic of dark}
We make the following observations about $\dark(D)$.
\begin{itemize}
\item By construction, $\dark(D) \in \Rook_+$.
\item Take $(r,c) \in D$.
If there are no $r' > r$ with $(r', c) \in \dark(D)$ 
and there are no $c' > c$ with $(r, c') \in \dark(D)$,
then $(r,c) \in \dark(D)$.
\end{itemize}
\end{rem}

Finally, we associate two statistics
to each diagram via its snow diagram.

\begin{defn}\label{defn.raj}
Let $D$ be a diagram.
The \definition{rajcode} of $D$,
$\rajcode(D)$, 
is the weak composition $\wt(\snow(D))$.
Let $\raj(D)$ denote $|\rajcode(D)|$,
the total number of cells in $\snow(D)$.
\end{defn}

\begin{exa}
Continuing with Example~\ref{exa.snow}, we have $\rajcode(D) = (3,3,2,1,2)$ and $\raj(D) = 11$.
\end{exa}

\begin{rem}
Recall that Pechenik, Speyer and Weigandt~\cite{PSW} 
define the statistics $\rajcode(\cdot)$ and $\raj(\cdot)$ on permutations
using increasing subsequences. 
We show that our $\rajcode$ and $\raj$ on Rothe diagrams 
agree with their definitions in Theorem~\ref{T: equivalence of rajcode on perm}.
Therefore, our construction on Rothe diagrams is a diagrammatic way
to compute the leading monomial and degree of $\topGro_w$.
In addition, we notice that positions of dark clouds in $\snow(RD(w))$
are connected to the Schensted insertion and Viennot's geometric construction.
These connections are explored in \S\ref{S: permutations}.
\end{rem}


\section{Proof of Theorem~\ref{T: Top Las}}
\label{S: weak composition}

To prove Theorem~\ref{T: Top Las},
we study top Lascoux polynomials via snow diagrams of key diagrams.
With a slight abuse of notation, we
define $\rajcode(\alpha) := \rajcode(D(\alpha))$, $\raj(\alpha) := \raj(D(\alpha))$
and $\dark(\alpha) = \dark(D(\alpha))$
for $\alpha \in C_+$.
We start by introducing some definitions.

\begin{defn}
A weak composition $\alpha$ is called \definition{snowy}
if its positive entries are all distinct. 
\end{defn}

Our main goal in this section is to establish 
Theorem~\ref{T: Top Las}:

\begin{manualtheorem}{\ref{T: Top Las}}
Let $\alpha$ and $\gamma$ be two weak compositions. 
\begin{enumerate}
\item[(a)]
The polynomial $\topLas_\alpha$ has leading monomial $x^{\rajcode(\alpha)}$.
\item[(b)] We have
$\topLas_\alpha$ is a scalar multiple of $\topLas_{\gamma}$
if and only if $\rajcode(\alpha) = \rajcode(\gamma)$.
\item[(c)] 
We say $\alpha$ is snowy if its
positive entries are distinct.
If $\alpha$ is snowy, then $x^{\rajcode(\alpha)}$ has coefficient 1 in $\topLas_\alpha$.
Moreover, 
there exists exactly one snowy weak composition $\gamma'$ 
such that $\rajcode(\gamma) = \rajcode(\gamma')$.
\end{enumerate}  
\end{manualtheorem}

This task is broken into four major lemmas
established in the following four subsections. 
In Subsection~\ref{SS: Top Las: Exist}, 
we use $K$-Kohnert diagrams to establish the first major lemma:

\begin{lemma}
\label{L: Top Las 1}
The polynomial $\fLb_\alpha$ has the term $x^{\rajcode(\alpha)} \beta^{\raj(\alpha) - |\alpha|}$.  
\end{lemma}

 


Lemma~\ref{L: Top Las 1} proves
$\topLas_\alpha$ has degree at least $\raj(\alpha)$.
To show $\topLas_\alpha$ indeed has degree $\raj(\alpha)$,
we need the following equivalence relation on weak compositions. 

\begin{defn}
Let $\alpha$ and $\gamma$ be two weak compositions. 
We say $\alpha$ is \definition{rajcode equivalent} to $\gamma$, 
denoted as $\alpha \sim \gamma$, 
if $\rajcode(\alpha) = \rajcode(\gamma)$. 
\end{defn}

\begin{exa}
\label{E: sim of weak comp}
Let $\alpha = (2,0,4,3,1)$ and $\gamma = (3,1,4,3,1)$.
Then we have:
\begin{align*}
D(\alpha) = 
\raisebox{0.95cm}{
\begin{ytableau}
\none[1] & \cdot& \cdot & \none & \none \cr
\none[2] & \none& \none & \none & \none \cr
\none[3] & \cdot& \cdot & \cdot & \cdot \cr
\none[4] & \cdot& \cdot & \cdot & \none \cr
\none[5] & \cdot& \none & \none & \none \cr
\end{ytableau}}\quad,
\quad \quad
\snow(D(\alpha)) = 
\raisebox{0.95cm}{
\begin{ytableau}
\none[1] & \cdot& \bullet & \snowflake & \snowflake \cr
\none[2] & \snowflake &  & \snowflake & \snowflake \cr
\none[3] & \cdot& \cdot & \cdot & \bullet \cr
\none[4] & \cdot& \cdot & \bullet & \none \cr
\none[5] & \bullet& \none & \none & \none \cr
\end{ytableau}}\quad,
\\
\end{align*}

\begin{align*}
D(\gamma) = 
\raisebox{0.95cm}{
\begin{ytableau}
\none[1] & \cdot& \cdot & \cdot & \none \cr
\none[2] & \cdot& \none & \none & \none \cr
\none[3] & \cdot& \cdot & \cdot & \cdot \cr
\none[4] & \cdot& \cdot & \cdot & \none \cr
\none[5] & \cdot& \none & \none & \none \cr
\end{ytableau}}\quad,
\quad \quad
\snow(D(\gamma)) = 
\raisebox{0.95cm}{
\begin{ytableau}
\none[1] & \cdot& \bullet & \cdot & \snowflake \cr
\none[2] & \cdot &  & \snowflake & \snowflake \cr
\none[3] & \cdot& \cdot & \cdot & \bullet \cr
\none[4] & \cdot& \cdot & \bullet & \none \cr
\none[5] & \bullet& \none & \none & \none \cr
\end{ytableau}}\quad.
\end{align*}

Be aware that the cell $(2,2)$ is not in
$\snow(D(\alpha))$ or $\snow(D(\gamma))$.
Observe that $\rajcode(\alpha) = (4,3,4,3,1) = \rajcode(\gamma)$,
so $\alpha \sim \gamma$.
\end{exa}

In Subsection~\ref{SS: equivalence},
we study this equivalence relation.
We show that snowy weak compositions form a complete set of representatives:

\begin{lemma}
\label{L: snowy weak composition in equivalence class}
For each equivalence class of $\sim$, 
there is a unique $\alpha$ such that $\alpha$ is snowy. 
Moreover, if $\gamma \sim \alpha$ and $\alpha$ is snowy, 
then $\gamma_r \geqslant \alpha_r$ for all $r$.
In other words, a snowy weak composition is the
unique entry-wise minimum in each equivalence class. 
\end{lemma}


In Subsection~\ref{SS: snowy}, 
we focus on $\topLas_\alpha$ for snowy $\alpha$ and give a recursive description of~$\topLas_\alpha$,
which leads to the third major lemma.

\begin{lemma}
\label{L: snowy leading}
If $\alpha$ is snowy, 
then $x^{\rajcode(\alpha)}$ is the leading monomial of 
$\topLas_\alpha$ with coefficient 1.
\end{lemma}
Finally, we devote the Subsection~\ref{SS: last major lemma} to proving the last
major lemma:

\begin{lemma}
\label{L: Top Las similar}
If $\alpha \sim \gamma$, 
then $\topLas_\alpha = c\topLas_\gamma$ for some $c \neq 0$.
\end{lemma}

Once we have these four major lemmas, 
we can easily check Theorem~\ref{T: Top Las}.


\begin{proof}
First, statement (c) follows from Lemma~\ref{L: snowy weak composition in equivalence class} and Lemma~\ref{L: snowy leading}. 

Given a weak composition $\alpha$. Let $\beta$ be the unique snowy weak composition such that $\alpha \sim \beta$. Statement (a) follows from  Lemma~\ref{L: snowy leading} and Lemma~\ref{L: Top Las similar}.

For statement (b), the backward direction is just Lemma~\ref{L: Top Las similar}.
For the forward direction, 
if $\topLas_\alpha$ is a scalar multiple of $\topLas_\gamma$,
then they have the same leading monomial.
By statement (a), we have $\rajcode(\alpha) = \rajcode(\gamma)$.
\end{proof}


\subsection{Proof of Lemma~\ref{L: Top Las 1}}
\label{SS: Top Las: Exist}

We show the monomial $ x^{\rajcode(\alpha)}\beta^{\raj(\alpha) - |\alpha|}$ 
exists in $\fLb_\alpha$.
We give an algorithm whose output is a $K$-Kohnert diagram for $\alpha$, which has the same underlying diagram as $\snow(D(\alpha))$. First, observe that $\snow(D(\alpha))$ contains no dark clouds if and only if $\alpha$ contains only zero entries. In this case, $\topLas_\alpha = 1$ and $\rajcode(\alpha)$ only has zero entries. Our claim is immediate. 
In the rest of this subsection, we assume $\alpha$ is a weak composition
with at least one positive entry, and thus $\snow(D(\alpha))$ has at least one dark cloud.
To describe the algorithm, 
we introduce two useful moves on ghost diagrams.
\begin{defn}
Let $D$ be a ghost diagram.
Let $(r,c)$ be a non-ghost cell in $D$
and let $(r', c)$ be the highest empty space in column $c$.
If $r' < r$,
let $UP_{(r,c)}(D)$ be the diagram we get after 
moving $(r,c)$ to $(r',c)$. 
Let $UP^G_{(r,c)}(D)$ be the diagram we get after 
moving $(r,c)$ to $(r',c)$
and putting a ghost on $(r,c)$ and all empty spaces between $(r,c)$ and $(r',c)$.
If $r' > r$, define 
$UP^G_{(r,c)}(D) = UP_{(r,c)}(D) = D$.
\end{defn}

\begin{rem}
\label{R: K-Kohnert move}
Assume $UP_{(r,c)}$ or $UP^G_{(r,c)}$ moves a cell to $(r',c)$.
Then this move can be achieved by a sequence of $K$-Kohnert moves
if both of the following conditions hold for each $r' < j \leq r$:
\begin{itemize}
\item If $(j, c)\notin D$,
then $D$ has no cell to the right of column $c$ in row $j$.
\item If $(j, c)\in D$, then it is not a ghost cell.
\end{itemize}
\end{rem}


Now we can describe the algorithm.
Let $D^{0} = D(\alpha)$. Recall by Remark~\ref{R: Basic of dark}, there is at most one dark cloud in each column of $\snow(D(\alpha))$. 
We can label all the dark clouds as $(r_1,c_1),\dots, (r_{m},c_{m})$ where $c_1<c_2\dots <c_m$ for some $m\geqslant 1$. We iterate $i$ from $1$ to $m$.
At iteration $i$, compute

\begin{align}
\label{EQ: UP}
D^i = UP^G_{(r_i,c_i)} \circ UP_{(r_i,c_i + 1)} \cdots \circ UP_{(r_i,\alpha_{r_i})} (D^{i-1})\,.    
\end{align}

\begin{exa}
Consider $\alpha = (1,3,4,0,4,3)$, we compute its snow diagram
and we have the dark clouds at $(2,1),(3,2),(6,3),(5,4)$. We compute $D^4$ according to the above algorithm. 
\[
\snow(D(\alpha)) = 
\raisebox{1cm}{\begin{ytableau}
\none[1] & \cdot & \snowflake & \snowflake & \snowflake \cr
\none[2] & \bullet & \cdot & \cdot & \snowflake \cr
\none[3] & \cdot & \bullet & \cdot & \cdot \cr
\none[4] & \none & \none & \snowflake & \snowflake\cr
\none[5] & \cdot& \cdot & \cdot & \bullet \cr
\none[6] & \cdot& \cdot & \bullet \cr
\end{ytableau}}\quad
D^0 = 
\raisebox{1cm}{\begin{ytableau}
\none[1] & \cdot \cr
\none[2] & \cdot & \cdot & \cdot \cr
\none[3] & \cdot & \cdot & \cdot & \cdot \cr
\none[4] & \none \cr
\none[5] & \cdot & \cdot & \cdot & \cdot \cr
\none[6] & \cdot & \cdot & \cdot \cr
\end{ytableau}}\, \xrightarrow[(2,1)]{}
\raisebox{1cm}{\begin{ytableau}
\none[1] & \cdot & \cdot & \cdot \cr
\none[2] & \cdot \cr
\none[3] & \cdot & \cdot & \cdot & \cdot \cr
\none[4] & \none \cr
\none[5] & \cdot & \cdot & \cdot & \cdot \cr
\none[6] & \cdot & \cdot & \cdot \cr
\end{ytableau}}\,
\]
\[
 \xrightarrow[(3,2)]{}\raisebox{1cm}{\begin{ytableau}
\none[1] & \cdot & \cdot & \cdot & \cdot \cr
\none[2] & \cdot & \cdot & \cdot \cr
\none[3] & \cdot & \X  \cr
\none[4] & \none \cr
\none[5] & \cdot & \cdot & \cdot & \cdot \cr
\none[6] & \cdot & \cdot & \cdot \cr
\end{ytableau}}\,\xrightarrow[(6,3)]{}
\raisebox{1cm}{\begin{ytableau}
\none[1] & \cdot & \cdot & \cdot & \cdot \cr
\none[2] & \cdot & \cdot & \cdot \cr
\none[3] & \cdot & \X &  \cdot \cr
\none[4] & \none & \none & \X \cr
\none[5] & \cdot & \cdot & \cdot & \cdot \cr
\none[6] & \cdot & \cdot & \X \cr
\end{ytableau}}\,\xrightarrow[(5,4)]{}
\raisebox{1cm}{\begin{ytableau}
\none[1] & \cdot & \cdot & \cdot & \cdot \cr
\none[2] & \cdot & \cdot & \cdot & \cdot \cr
\none[3] & \cdot & \X    &  \cdot    &  \X\cr
\none[4] & \none & \none & \X    & \X \cr
\none[5] & \cdot & \cdot & \cdot & \X \cr
\none[6] & \cdot & \cdot & \X \cr
\end{ytableau}} = D^{4}.
\]
\end{exa}

We observe that in the previous example, 
$D^4$ has the same underlying diagram as $\snow(D(\alpha))$.
This is true in general.

\begin{lem}
The labeled diagram $D^m$ defined by~(\ref{EQ: UP})
has the same underlying diagram as $\snow(D(\alpha))$.
\end{lem}
\begin{proof}
For a number $c$, 
we compare the column $c$ 
of $\snow(D(\alpha))$ and $D^m$.
If column $c$ of $\snow(D(\alpha))$ has no dark cloud,
then it is the same as column $c$ of $D(\alpha)$.
In this case, the algorithm will not move any cells 
in column $c$.
Thus, $D^m$ and $D(\alpha)$ also agree in column $c$.

Now suppose $\snow(D(\alpha))$ has a dark cloud
in column $c$, say at row $r$.
In the underlying diagram of $\snow(D(\alpha))$,
column $c$ is obtained from column $c$ of $D(\alpha)$ 
by filling all empty spaces above row $r$.  
On the other hand, consider what the algorithm does on column $c$.
It first might move cells above row $r$
and then it fills all empty spaces weakly above row $r$.
Thus, column $c$ in $D^m$ is the same as 
column $c$ of $\snow(D(\alpha))$ after ignoring the labels.  
\end{proof}

Next, we want to show $D^m$ produced by the algorithm
is in $\KKD(\alpha)$.
We just need to check each $UP_{(r,j)}$ and $UP^G_{(r,c)}$ in each iteration is a sequence of $K$-Kohnert moves. To that end, we first make the following observation about the diagram $D^i$.

\begin{lemma}
Let $c_0 = 0$. In $D^i$,
if a cell is strictly to the right of column $c_i$,
then there is a cell immediately on its left.
In other words, 
the diagram $D^i$ is left-justified if we ignore the first $c_i$ columns.
\end{lemma}

\begin{proof}
Prove by induction on $i$.
The lemma holds for $D^0$, which is left-justified.

Assume $D^{i-1}$ is left-justified if we ignore the first $c_{i-1}$ columns,
for some $i \geq 1$.
Consider an arbitrary cell $(r,c)$ in $D^i$
with $c > c_i$.
We show $(r, c-1)$ is in $D^i$ by considering two possibilities.
\begin{itemize}
\item[-] The cell $(r,c)$ is not in $D^{i-1}$.
Then during iteration $i$, a cell is moved to $(r,c)$,
which is the highest blank in column $c$ of $D^{i-1}$.
By our inductive hypothesis and $c-1 > c_{i-1}$, 
the highest blank in column $c-1$ of $D^{i-1}$
is weakly lower than row $r$.
Thus, $(r,c-1)$ is in $D^i$.
\item[-] Otherwise, $(r,c)$ is in $D^{i-1}$.
By our inductive hypothesis, 
$(r, c-1)$ is in $D^{i-1}$.
If $r \neq r_i$,
then we know that no cell from row $r$ is moved 
during iteration $i$.
Thus, $(r,c-1)$ is still in $D^i$.
If $r = r_i$,
then there are no empty spaces above $(r,c)$
in $D^{i-1}$.
By our inductive hypothesis, 
there is no empty spaces above $(r,c-1)$,
so $(r,c-1)$ is still in $D^i$. \qedhere
\end{itemize}
\end{proof}

The above lemma shows that the diagram $D^i$ is left-justified if we ignore the first $c_i$ columns.
We will use this property to show that $D^m$ 
is in $\KKD(\alpha)$.
\begin{prop}\label{P: Top Las 1}
The above algorithm can be achieved by $K$-Kohnert moves,
so $D^m \in \KKD(\alpha)$.
\end{prop}

\begin{proof}
We focus on one iteration of the algorithm,
say iteration $i$.
We check the operators 
in~(\ref{EQ: UP}) can be achieved by $K$-Kohnert moves. 
We ignore all cells to the left of the column $c_i$ in $D^{i-1}$.
By the previous Lemma, this part of the diagram 
is left-justified.
The highest empty spaces in columns $c_i, \cdots, \alpha_{r_i}$
are going weakly up from left to right.
Moreover, the condition in Remark~\ref{R: K-Kohnert move} 
holds for all $(r_i, c_i), \cdots, (r_i, \alpha_{r_i})$.

Now $UP_{(r_i, \alpha_{r_i})}$
can be achieved by $K$-Kohnert moves.
After that, the conditions in Remark~\ref{R: K-Kohnert move} 
hold at each step for $(r_i, \alpha_{r_i} - 1), \dots, (r_i, c_i)$.
Following this logic, this iteration can be achieved by $K$-Kohnert moves. 
\end{proof}

Using Theorem~\ref{thm: k.kohnert}:

\begin{manuallemma}{\ref{L: Top Las 1}}
The polynomial $\fLb_\alpha$ has the term $x^{\rajcode(\alpha)} \beta^{\raj(\alpha) - |\alpha|}$.
\end{manuallemma} 

\subsection{Proof of Lemma~\ref{L: snowy weak composition in equivalence class}}
\label{SS: equivalence}

First, notice that we can recover the underlying diagram of $\snow(D(\alpha))$
from $\dark(\alpha)$.

\begin{lemma}
\label{L: dark recovers snow}
Let $\alpha$ be a weak composition. 
The underlying diagram of $\snow(D(\alpha))$
is:
\begin{align}
\label{Eq: snow recovered by dark}
\bigcup_{(r,c) \in \dark(\alpha)} ([r] \times \{c\}) \cup (\{r\} \times [c]).
\end{align}
\end{lemma}
\begin{proof}
First, we show that the elements of the set~(\ref{Eq: snow recovered by dark})
are cells in $\snow(D(\alpha))$.
Take $(r,c) \in \dark(\alpha)$.
We know $(r,c) \in D(\alpha)$.
Since $D(\alpha)$ is left-justified, 
$\{r\} \times [c] \subseteq D(\alpha)$.
Thus, these cells are in $\snow(D(\alpha))$.
By the construction of $\snow(D(\alpha))$,
the cells in~$[r] \times \{c\}$ are also 
in $\snow(D(\alpha))$.

Now suppose there is a cell $(r, c)$ 
in $\snow(D(\alpha))$
that is not in the 
set~(\ref{Eq: snow recovered by dark}).
Then there is no $r' > r$ with $(r', c) \in \dark(D)$,
which implies $(r,c)$ is not a snowflake 
in $\snow(D(\alpha))$.
Thus, $(r,c) \in D(\alpha)$.
Also, there is no $c' > c$ with $(r, c') \in \dark(D)$.
By Remark~\ref{R: Basic of dark},
$(r,c) \in \dark(D)$.
Thus, $(r, c)$ is in the 
set~(\ref{Eq: snow recovered by dark}),
which is a contradiction.
\end{proof}

Furthermore, 
we can recover $\dark(\alpha)$ from $\rajcode(\alpha)$.

\begin{lemma}
\label{L: same rajcode => same dark}
Let $\alpha, \gamma$ be weak compositions.
If $\rajcode(\alpha) = \rajcode(\gamma)$,
then $\dark(\alpha) = \dark(\gamma)$.
\end{lemma}
\begin{proof}
We prove the two diagrams $\dark(\alpha)$ 
and $\dark(\gamma)$ agree on each row $r$,
by a reverse induction on $r$.
The base case is immediate.
Suppose $r$ is large enough
such that $\alpha_i = \gamma_i = 0$
if $i > r$.
Then $\dark(\alpha)$ and $\dark(\gamma)$
clearly agree on row $r$ and underneath.

Next, we show that the value $\rajcode(\alpha)_r$ and cells in $\dark(\alpha)$ under row $r$ 
determines whether $\dark(\alpha)$ has a cell on row $r$.
Moreover, if such a cell exists, its column index is also determined. 

Let $r\geqslant 1$.
Define
$$B_r :=\{c: \textrm{There are no dark clouds under } (r,c) \textrm{ in } \snow(D(\alpha) \}.$$

The complement of $B_r$ is $\overline{B_r} := \mathbb{Z}_{>0} - B_r = \{c: (r',c) \in \dark(\alpha) \textrm{ for some } r' > r \}$.
For $c \in \overline{B_r}$,
$(r,c)$ of $\snow(D(\alpha))$ is a snowflake or an unlabeled cell.
If there is no dark cloud on row $r$ of $\snow(D(\alpha))$,
$\rajcode(\alpha)_r = |\overline{B_r}|$.
Otherwise, we assume the dark cloud is at $(r,c)$ for some $c \in B_r$.
Then row $r$ of $\snow(D(\alpha))$ has cells on 
$(r,c')$ for $c' \in \overline{B_r}$ or $c' \leq c$.
Suppose $c$ is the $i$\textsuperscript{th} smallest number in $B_r$.
We have $\rajcode(\alpha)_r = i + |\overline{B_r}|$.

Consequently, $\rajcode(\alpha)_r$ and $\dark(\alpha)$ under row $r$
uniquely determines row $r$ of $\dark(\alpha)$.
If we assume $\dark(\alpha)$ and $\dark(\gamma)$
agree underneath row $r$ as our inductive hypothesis,
then they also agree on row $r$ since $\rajcode(\alpha)_r = \rajcode(\gamma)_r$.
The induction is finished. 
\end{proof}

Now we have two equivalent ways of describing rajcode equivalence.
\begin{prop}
\label{P: equivalence}
Let $\alpha$ and $\gamma$ be two weak compositions. 
The following are equivalent:
\begin{enumerate}
\item $\alpha \sim \gamma$;
\item $\dark(\alpha) = \dark(\gamma)$.
\item The underlying diagrams of $\snow(D(\alpha))$ and  $\snow(D(\gamma))$
are the same;
\end{enumerate}
\end{prop}
\begin{proof}
By Lemma~\ref{L: same rajcode => same dark}, (1) implies (2).
By Lemma~\ref{L: dark recovers snow}, (2) implies (3).
Clearly, (3) implies (1).
\end{proof}

Our next goal is to find representatives
of rajcode equivalence classes. 
At the end of this subsection, 
we will see snowy weak compositions form
a complete set of representatives. 
To understand snowy weak compositions, 
we start with the following observation. 

\begin{rem}
\label{R: snowy}
For a weak composition $\alpha$,
the following are equivalent:
\begin{itemize}
\item $\alpha$ is snowy.
\item The rightmost cell in each row of $D(\alpha)$
are in different columns. 
\item The rightmost cell in each row of $D(\alpha)$
is a dark cloud in $\snow(D(\alpha))$. 
\end{itemize}
\end{rem}

One advantage of working with snowy weak compositions
is that 
we can tell their 
$\rajcode(\cdot)$ and $\raj(\cdot)$ easily:
\begin{lemma}
\label{L: dark, rajcode and raj of snowy weak composition}
Let $\alpha$ be a snowy weak composition.
Then the following statements hold.
\begin{enumerate}
\item $\dark(\alpha) = \{(r,\alpha_r): \alpha_r > 0\}$,\label{snowy.weak.1}
\item $\rajcode(\alpha)_r = \alpha_r + |\{r' > r: \alpha_r < \alpha_r'\}|$, and \label{snowy.weak.2}
\item $\raj(\alpha) = \sum_r (\alpha_r + |\{(r, r'): \alpha_r < \alpha_r', r < r'\}|) = |\alpha| + |\{(r, r'): r < r', \alpha_r < \alpha_{r'}\}|$.\label{snowy.weak.3}
\end{enumerate}
\end{lemma}
\begin{proof}
(1) follows from Remark~\ref{R: snowy}.
(2) follows from (1) and Lemma~\ref{L: dark recovers snow}, and 
(3) immediately follows from (2). 
\end{proof}

As a consequence, we have the following rule which tells us 
how $\rajcode(s_i\alpha)$ differs from 
$\rajcode(\alpha)$ when $\alpha$ is snowy. 

\begin{cor}
\label{C: snowy changes}
Let $\alpha$ be a snowy weak composition and consider $i$ with $\alpha_i > \alpha_{i+1}$.
Then $\rajcode(s_i\alpha) = s_i\rajcode(\alpha) + e_i$,
where $e_i$ is the weak composition with
1 on its $i$\textsuperscript{th} entry and 0 elsewhere.
\end{cor}

The second advantage of working with snowy weak compositions is that
they are in bijection with $\Rook_+$.

\begin{lem}
\label{L: snowy biject rook}
The map $\dark(\cdot)$ is a bijection
from $\{ \alpha \in C_+: \alpha \textrm{ is snowy}\}$
to $\Rook_+$.
Its inverse $\dark^{-1}(\cdot)$ 
is given by $\dark^{-1}(R) = \alpha$
where \begin{align*}
\alpha_r =
\begin{cases}
0 &\text{if row $r$ of $R$ is empty;} \\
c &\text{if $(r,c) \in R$.} 
\end{cases}
\end{align*}
\end{lem}
\begin{proof}
Follows from Remark \ref{R: snowy}.
\end{proof}

We are ready to show that they are representatives of all equivalence classes.

\begin{manuallemma}{\ref{L: snowy weak composition in equivalence class}}
For each equivalence class of $\sim$, 
there is a unique $\alpha$ such that $\alpha$ is snowy. 
Moreover, if $\gamma \sim \alpha$ and $\alpha$ is snowy, 
then $\gamma_r \geqslant \alpha_r$ for all $r$.
In other words, a snowy weak composition is the
unique entry-wise minimum in each equivalence class. 
\end{manuallemma}
\begin{proof}
Let $\gamma$ be an arbitrary weak composition. 
First, we construct a snowy $\alpha$ 
such that $\alpha \sim \gamma$.
We know $\dark(\gamma) \in \Rook_+$.
We send it to a snowy $\alpha$
using the map in Lemma~\ref{L: snowy biject rook}.
Then $\dark(\alpha) = \dark(\gamma)$.
By Proposition~\ref{P: equivalence},
$\alpha \sim \gamma$.

Next, take a positive integer $r$.
If $\alpha_r = 0$, then $\gamma_r \geqslant \alpha_r$ trivially.
Otherwise, we know $(r, \alpha_r) \in \dark(\alpha) = \dark(\gamma)$.
Thus, $\gamma_r \geqslant \alpha_r$.

Finally, we establish the uniqueness of this snowy $\alpha$.
Assume $\alpha'$ is a snowy weak composition
such that $\alpha' \sim \gamma$.
Then $\alpha'_r \geqslant \alpha_r$ and $\alpha_r \geqslant \alpha'_r$
for all $r \in \Z_{>0}$, so $\alpha = \alpha'$.
\end{proof}

A snowy weak composition has more snowflakes in its snow diagram 
than any others in its equivalence class; hence the name.
Say $\alpha \sim \gamma$ and
$\alpha$ is snowy while $\gamma$ is not. 
By Lemma~\ref{L: snowy weak composition in equivalence class}, $|\alpha| <  |\gamma|$.
On the other hand, 
the number of snowflakes in $\snow(D(\alpha))$ 
(resp. $\snow(D(\gamma))$)
is $\raj(\alpha) - |\alpha|$ (resp. $\raj(\gamma) - |\gamma|$). 
Since $\raj(\alpha) = \raj(\gamma)$, $\snow(D(\alpha))$
has more snowflakes than $\snow(D(\gamma))$.

\subsection{Proof of Lemma~\ref{L: snowy leading}}
\label{SS: snowy}

By Lemma~\ref{L: Top Las 1},
$\topLas_\alpha$ has degree at least $\raj(\alpha)$.
Next, we can show the degree of $\topLas_\alpha$
equals to $\raj(\alpha)$ when $\alpha$ is snowy. 

\begin{lemma}
\label{L: Snowy degree}
Let $\alpha$ be a snowy weak composition.
The $\beta$-degree of $\fLb_\alpha$ 
is $\raj(\alpha) - |\alpha|$,
so the degree of $\topLas_\alpha$ is $\raj(\alpha)$.
\end{lemma}
\begin{proof}
We prove the result by induction on 
\begin{align*}
\ell(\alpha) := |\{(i,j) \mid\, \alpha_i < \alpha_j \text{ and } i < j\}|.
\end{align*}
For the base case, if $\ell(\alpha) = 0$,
then $\alpha$ is weakly decreasing.
The polynomial $\fLb_\alpha$ is an monomial with $\beta$-degree $0$.
Correspondingly, $\raj(\alpha) = |\alpha|$.

Now if $\ell(\alpha) > 0$,
we can find $i$ with $\alpha_i < \alpha_{i+1}$.
By Corollary~\ref{C: snowy changes},
$\raj(s_i \alpha) = \raj(\alpha) - 1$.
Notice that $\ell(s_i \alpha) = \ell(\alpha) - 1$.
By our inductive hypothesis, 
the $\beta$-degree of $\fLb_{s_i \alpha}$
is $\raj(s_i \alpha) - |\alpha| = \raj(\alpha) - 1 - |\alpha|$.
By the recursive definition of Lascoux polynomials,
$$
\fLb_\alpha 
= \pi_i(\fLb_{s_i\alpha}) + \beta \pi_i(x_{i+1}\fLb_{s_i\alpha}).
$$
The $\beta$-degree in $\fLb_\alpha$
is at most $\raj(\alpha) - |\alpha|$.
Lemma~\ref{L: Top Las 1} implies the 
$\beta$-degree of $\fLb_\alpha$ is at least 
$\raj(\alpha) - |\alpha|$,
so the inductive step is finished.
\end{proof}


Combine with Lemma~\ref{L: dark, rajcode and raj of snowy weak composition},
we have: 
\begin{cor}
Let $\alpha$ be a snowy weak composition.
The degree of $\topLas_\alpha$
is $|\alpha| + |\{(r, r'): r < r', \alpha_r < \alpha_{r'}\}|$.
\end{cor}

Now we can describe $\topLas_\alpha$ for snowy 
$\alpha$ recursively.
\begin{lemma}\label{L: recursive top Las}
Let $\alpha$ be a snowy weak composition. 
Then
\begin{align}
\topLas_\alpha = \begin{cases}
x^\alpha & \text{if $\alpha_1 \geqslant \alpha_2 \geqslant \cdots$} \\
\pi_i (x_{i+1} \topLas_{s_i\alpha}) &\text{if $\alpha_i<\alpha_{i+1}$.}
\end{cases}
\end{align}
\end{lemma}
\begin{proof}
When $\alpha$ is weakly decreasing, 
our rule is immediate.
Now assume $\alpha_i < \alpha_{i+1}$ for some $i\in \Z_{>0}$.
By Corollary~\ref{C: snowy changes},
$\raj(s_i \alpha) = \raj(\alpha) - 1$.
We write $\fLb_{s_i \alpha}$
as $g + \beta^{\raj(\alpha) - 1 - |\alpha|} \: \topLas_{s_i \alpha}$
for some $g \in \Z[x_1, x_2, \cdots][\beta]$
with $\beta$-degree less than 
$\raj(\alpha) - 1 - |\alpha|$.
Now we write $\fLb_\alpha$ as
\begin{align*}
\fLb_\alpha & = \pi_i(\fLb_{s_i \alpha}) 
+ \beta \pi_i(x_{i+1}\fLb_{s_i \alpha})  \\ 
& =\pi_i(\fLb_{s_i \alpha})
+ \beta \pi_i(x_{i+1}g) + \beta^{\raj(\alpha) - |\alpha|} \pi_i(x_{i+1}\topLas_{s_i \alpha})   
\end{align*}
When we extract the coefficient of $\beta^{\raj(\alpha) - |\alpha|}$,
the left-hand side is $\topLas_\alpha$.
On the right-hand side, 
the first two terms are ignored and 
we get $\pi_i(x_{i+1}\topLas_{s_i \alpha})$.
\end{proof}

Combining Lemma~\ref{L: Top Las 1} and Lemma~\ref{L: Snowy degree},
we know $x^{\rajcode(\alpha)}$ appears in $\topLas_\alpha$
when $\alpha$ is snowy. 
Next, we show this monomial
is the leading monomial of $\topLas_\alpha$.
We start with the following observation 
about the operator $f \mapsto \pi_i(x_{i+1} f)$.
\begin{rem}
\label{R: leading}
Let $\gamma$ be a monomial.
We may describe the leading monomial of
$\pi_i(x_{i+1}x^\gamma)$ and its coefficient as follows. 
\begin{itemize}
\item If $\gamma_i > \gamma_{i+1}$,
then $x_i x^{s_i \gamma}$ is the leading monomial
with coefficient 1.
\item If $\gamma_i = \gamma_{i+1}$,
then $\pi_i(x_{i+1}x^\gamma) = 0$. 
\item If $\gamma_i < \gamma_{i+1}$,
then $x_i x^\gamma$ is the leading monomial 
with coefficient $-1$. 
\end{itemize}
\end{rem}

We can understand how the operator
$f \mapsto \pi_i(x_{i+1} f)$ changes 
the leading monomial of polynomial $f$
satisfying certain conditions. 

\begin{lemma}
\label{L: monomials from operator}
Take $f \in \Z[x_1, x_2, \cdots ]$ with $f \neq 0$.
Assume $x^\alpha$ is the leading monomial in $f$ 
with coefficient $c \neq 0$.
Pick an $i\in \Z_{>0}$ such that $\alpha_{i} > \alpha_{i+1}$.
Furthermore, assume for any monomial in $f$, 
its power of $x_{i}$ is at most $\alpha_{i}$.
Then $x_ix^{s_i \alpha}$ is the leading monomial in
$\pi_{i}(x_{i+1}f)$ with coefficient $c$. 
\end{lemma}
\begin{proof}
In this proof, we use ``$\geq$'' to denote the 
monomial order.
Let $\Gamma$ be the set of weak compositions $\gamma$
such that $x^\gamma$ appears in $f$.
Let $c_\gamma$ be the coefficient of $x^\gamma$ in $f$.
We may write $f = \sum_{\gamma \in \Gamma} c_\gamma x^\gamma$.
Then 
$\pi_{i}(x_{i+1}f) = \sum_{\gamma \in \Gamma} c_\gamma \pi_{i}(x_{i+1}x^\gamma)$.
By the remark above, $x_ix^{s_i \alpha}$ appears 
in $c_\alpha\pi_i(x_{i+1}x^\alpha)$ as the leading monomial
with coefficient $c_\alpha = c$.
It is enough to show the following claim.

\noindent\textbf{Claim:} 
Take $\gamma \in \Gamma$ such that 
$\pi_i(x_{i+1}x^{\gamma}) \neq 0$ 
(i.e. $\gamma_i \neq \gamma_{i+1}$). 
Let $x^{\gamma'}$ be the leading monomial 
in $\pi_i(x_{i+1}x^{\gamma})$.
If $x^{\gamma'}  \geq x_ix^{s_i \alpha}$,
then $\gamma = \alpha$.

\noindent\textbf{Proof:} 
Assume $\alpha \neq \gamma$.
Let $k$ be the largest index such that 
the power of $x_k$ differs in $x^{\gamma'}$ and $x_ix^{s_i \alpha}$.
By $x^{\gamma'}  \geq x_ix^{s_i \alpha}$,
the power of $x_k$ in $x^{\gamma'}$ is greater than
the power of $x_k$ in $x_ix^{s_i \alpha}$.
We must have $k \leqslant i+1$.
Otherwise,
$x^\gamma > x^\alpha$,
which contradicts $x^\alpha$ being the leading monomial in $f$.

Now we know 
$\gamma'$, $\alpha$ and $\gamma$ all agree after 
the $(i+1)^{th}$ entry.
Then
$\gamma_{i+1}'$ is at least
the power of $x_{i+1}$ in $x_ix^{s_i \alpha}$,
which is $\alpha_i$.
On the other hand, by $x^\gamma \leqslant x^\alpha$,
$\gamma_{i+1} \leqslant \alpha_{i+1}$.
Thus, 
\begin{equation}
\label{EQ: chain inequality}
\gamma_{i+1} \leq  \alpha_{i+1}<\alpha_i \leq \gamma_{i+1}'.
\end{equation}
If $\gamma_{i} < \gamma_{i+1}$,
Remark~\ref{R: leading} implies $\gamma_{i+1}'= \gamma_{i+1}$,
which is impossible.
Thus, we must have $\gamma_i > \gamma_{i+1}$. 
By Remark~\ref{R: leading} again, $\gamma_{i+1}' = \gamma_i$.
By the assumptions in the statement of the lemma, $\gamma_i \leqslant \alpha_i$,
so $\gamma_{i+1}' = \gamma_i = \alpha_i$.

Next, 
$\gamma_{i}'$ is at least
the power of $x_{i}$ in $x_ix^{s_i \alpha}$,
which is $\alpha_{i+1} + 1$.
Remark~\ref{R: leading} implies $\gamma_i' = \gamma_{i+1} + 1$.
Thus, $\gamma_{i+1} \geqslant \alpha_{i+1}$.
By~(\ref{EQ: chain inequality}), $\gamma_{i+1} = \alpha_{i+1}$.

Now we know $k < i$
and $\gamma_j = \alpha_j$ for $j = i$ or $i+1$.
Thus, $\gamma_j = \alpha_j$ for all $j > k$,
so~$x^\gamma > x^\alpha$, which is a contradiction. 
\end{proof}

Now we can establish our third major lemma. 
\begin{manuallemma}{\ref{L: snowy leading}}
If $\alpha$ is snowy, 
then $x^{\rajcode(\alpha)}$ is the leading monomial of 
$\topLas_\alpha$ with coefficient 1.
\end{manuallemma}
\begin{proof}
We prove the result by induction on 
\begin{align*}
\ell(\alpha) := |\{(i,j) \mid\, \alpha_i < \alpha_j \text{ and } i < j\}|.
\end{align*}

If $\ell(\alpha)=0$, then $\alpha$ is weakly decreasing, 
then $\fLb_\alpha = 
x^\alpha = x^{\rajcode(\alpha)}$.
Our claim is immediate. 

Now if $\ell(\alpha) > 0$,
we can find $r$ with $\alpha_r < \alpha_{r+1}$.
Pick the largest such $r$.
For our inductive hypothesis,
assume $x^{\rajcode(s_r \alpha)}$ is the leading monomial 
of $\topLas_{s_r \alpha}$ with coefficient 1.

By the maximality of $r$, 
$\alpha_{r+1} \geqslant \alpha_{r+2} 
\geqslant \alpha_{r+3} \geqslant \cdots$.
Thus, in any $K$-Kohnert diagram of $s_r \alpha$,
there cannot be more than $\alpha_{r+1}$ cells
in row $r$.
In other words, 
for any monomial of $\topLas_{s_r \alpha}$,
the power of $x_r$ is at most $\alpha_{r+1}$.
Lemma~\ref{L: monomials from operator} implies that $x_r x^{s_r \rajcode(s_r \alpha)}$
is the leading monomial of $\topLas_{\alpha}$ with coefficient 1. 
Finally, by Corollary~\ref{C: snowy changes},
$x_r x^{s_r \rajcode(s_r \alpha)} = x^{\rajcode(\alpha)}$.
\end{proof}


\subsection{Proof of Lemma~\ref{L: Top Las similar}}
\label{SS: last major lemma}
We first derive two consequences 
of $\alpha \sim \gamma$.
We start with the following definition.
\begin{defn}
Let $D$ be a diagram.
Let $\overline{D} := \bigcup_{(r,c) \in D} [r] \times \{c\}$.
\end{defn}

In plain words, $\overline{D}$ is the diagram obtained
by filling the empty spaces above each cell of $D$.
Then $\overline{D(\alpha)}$ is completely determined 
by $\dark(\alpha)$:
\begin{lemma}
Let $\alpha$ be a weak composition.
Then $\overline{D(\alpha)} = \bigcup_{(r,c) \in \dark(\alpha)} [r] \times [c]$.
\end{lemma}
\begin{proof}
We show each side is a subset of the other. 
Take $(r_1,c_1) \in D(\alpha)$. 
By Remark~\ref{R: Basic of dark},
there is $(r_2, c_2) \in \dark(\alpha)$ such that 
$r_2 \geqslant r_1$ and $c_2 \geqslant c_1$.
Thus, $[r_1] \times \{c_1\} \subseteq [r_2] \times [c_2]$.

Take $(r_1,c_1) \in \dark(\alpha)$.
Thus, for any $c \in [c_1]$, 
$(r_1,c) \in D(\alpha)$.
Then $[r_1] \times \{c\} \subseteq \overline{D(\alpha)}$,
so $[r_1] \times [c_1] \subseteq \overline{D(\alpha)}$. 
\end{proof}

We have the following consequence of 
$\alpha \sim \gamma$.

\begin{cor}
\label{C: sim implies same overline}
If $\alpha \sim \gamma$,
then $\overline{D(\alpha)} = \overline{D(\gamma)}$.
\end{cor}
Notice that the converse is not true. 
If $\alpha = (1,2)$ and $\gamma = (0,2)$,
then $\overline{D(\alpha)} = [2] \times [2]
=  \overline{D(\gamma)}$.
However, $\alpha$ and $\gamma$ are not similar,
since $\dark(\alpha) = \{ (1,1), (2,2)\}$
and $\dark(\gamma) = \{(2,2)\}$.

Another nice consequence of $\alpha \sim \gamma$ one might expect
is $s_r \alpha \sim s_r \gamma$.
Unfortunately, this is not always true. 
It is easy to check $(0,1) \sim (1,1)$
but $s_1(0,1) = (1,0)$  and $s_1(1,1) = (1,1)$ are not similar. 
However, it is true
when $\alpha$ and $r$ satisfy 
the following condition.

\begin{lemma}
\label{L: s_i fix sim}
Let $\alpha$ be a weak composition and $r\in \Z_{>0}$.
Assume there exists $c$
such that $(r, c)\notin \snow(D(\alpha))$
but $(r+1, c)\in \snow(D(\alpha))$.
Then 
\begin{enumerate}
\item[(i)] $\alpha_{r+1} > \alpha_r$;
\item[(ii)] The diagram $\dark(s_r \alpha)$
is obtained from $\dark(\alpha)$
by switching row $r$ and row $r+1$;
\item[(iii)] For any $\gamma$ with
$\gamma \sim \alpha$,
we must have $\gamma_{r+1} > \gamma_{r}$ 
and~$s_r\alpha \sim s_r\gamma$.
\end{enumerate}
\end{lemma}

\begin{proof}
Since $(r,c)$ is not in $\snow(D(\alpha))$,
we can deduce two facts:
\begin{enumerate}
\item There are no dark clouds under row $r$ in column $c$, and 
\item $\alpha_r < c$.
\end{enumerate}

By (1), the cell $(r+1,c)$ in $\snow(D(\alpha))$ 
is not a dark cloud or a snowflake.
Thus, it is unlabeled and $(r+1,c) \in D(\alpha)$.
By Remark~\ref{R: Basic of dark},
there must be a $c' > c$
such that $(r+1, c')$ is a dark cloud 
in $\snow(D(\alpha))$.
This implies $\alpha_{r+1} > c$.
By (2), 
we have $\alpha_{r+1} > \alpha_r$, 
proving (i).
Also by (2), 
the dark cloud in row $r$ of 
$\snow(D(\alpha))$,
if exists, 
is in the first $c-1$ columns. 
Thus, $\dark(s_r\alpha)$ is obtained
from $\dark(\alpha)$ by switching 
row $r$ and row $r+1$, proving (ii).

Now consider any  $\gamma \sim \alpha$.
By Proposition~\ref{P: equivalence},
$\snow(D(\gamma))$ and $\snow(D(\alpha))$
have the same underlying diagram.
By (ii), 
$\dark(s_r\gamma)$ is obtained
from $\dark(\gamma)$ by switching 
row $r$ and row $r+1$.
Since $\dark(\alpha) = \dark(\gamma)$,
we have $\dark(s_r\alpha) = \dark(s_r\gamma)$,
so $s_r \alpha \sim s_r \gamma$.
\end{proof}


These two consequences of $\alpha \sim \gamma$
allow us to prove the last main Lemma. 

\begin{manuallemma}{\ref{L: Top Las similar}}
If $\alpha \sim \gamma$, 
then $\topLas_\alpha = c\topLas_\gamma$ for some $c \neq 0$.
\end{manuallemma}
\begin{proof}
By Lemma~\ref{L: snowy weak composition in equivalence class} 
it is enough to assume $\gamma$ is snowy, and we proceed by induction on $\raj(\alpha)$.
The base case is $\raj(\alpha) = 0$,
which implies $\alpha$ only has 0s. 
Our claim is immediate. 

Now assume $\raj(\alpha) > 0$.
Consider the diagram $\overline{D(\alpha)}$.
Clearly, the underlying diagram of any $K$-Kohnert diagram
of $\alpha$ will be a subset of $\overline{D(\alpha)}$.
In other words, any monomial in $\topLas_{\alpha}$
must divide $x^{\wt(\overline{D(\alpha)})}$.

If the underlying diagram of $\snow(D(\alpha))$
is $\overline{D(\alpha)}$, 
then $x^{\wt(\overline{D(\alpha)})}$ is the only monomial
in $\topLas_{\alpha}$.
On the other hand, 
Corollary~\ref{C: sim implies same overline} gives
$\overline{D(\alpha)} = \overline{D(\gamma)}$.
By the same argument, $x^{\wt(\overline{D(\alpha)})}$ is the only monomial
in $\topLas_{\gamma}$.
Our claim holds. 

Otherwise, we can find $(r,c) \in \overline{D(\alpha)}$
but not in $\snow(D(\alpha))$.
Choose the $(r,c)$ with the largest $r$.
First, we know $(r,c) \notin D(\alpha)$,
which implies $(r+1, c) \in \overline{D(\alpha)}$.
By the maximality of $r$,
$(r+1, c)$ is in $\snow(D(\alpha))$.
We invoke Lemma~\ref{L: s_i fix sim}
and conclude $\alpha_{r+1} > \alpha_{r}$,
$\gamma_{r+1} > \gamma_{r}$
and $s_i \alpha \sim s_i \gamma$.
Since $\gamma$ is snowy,
by Corollary~\ref{C: snowy changes},
we know $\raj(s_r \gamma) = \raj(\gamma) - 1$,
which implies $\raj(s_r \alpha) = \raj(\alpha) - 1$.
By our inductive hypothesis, 
$\topLas_{s_r \alpha} = c \topLas_{s_r \gamma}$
for some $c \neq 0$.

We may write 
$\fLb_{s_r \alpha}$ 
as $\beta^{\raj(s_r \alpha) - |\alpha|}\topLas_{s_r \alpha} + g$,
where $g$ has $\beta$-degree less than  
$\raj(s_r \alpha) - |\alpha|$.
Then 
\begin{align*}
\fLb_\alpha 
& = \pi_i(\fLb_{s_r \alpha}) + \beta \pi_i(x_{i+1}\fLb_{s_r \alpha})\\
& = \pi_i(\fLb_{s_r \alpha}) + \beta \pi_i(x_{i+1} g) + \beta^{\raj(\alpha) - |\alpha|} \pi_i(x_{i+1} \topLas_{s_r \alpha})    
\end{align*}

The first two terms on the right-hand side have $\beta$ degree 
less than $\raj(\alpha) - |\alpha|$.
Thus, the $\beta$-degree in $\fLb_\alpha$ is at most $\raj(\alpha) - |\alpha|$.
By Lemma~\ref{L: Top Las 1},
the $\beta$-degree in $\fLb_\alpha$ is $\raj(\alpha) - |\alpha|$.
Extract the coefficient of $\beta^{\raj(\alpha) - |\alpha|}$ and get
$$
\topLas_\alpha = \pi_i(x_{i+1} \topLas_{s_r \alpha}) = c\pi_i(x_{i+1} \topLas_{s_r \gamma}) = c\topLas_{\gamma},
$$
by Lemma~\ref{L: recursive top Las}.
\end{proof}

